% Standalone section: requires standard LaTeX, amsmath and amssymb.
% Input from the parent research report with \input{interpreter/interpreter_results}.
\section{Interpreter compression: a local saving and two integration barriers}
\label{sec:interpreter-compression}

The interpreter investigation produced a new local unary-register increment
routine that permits deletion of one working state, together with explicit
reasons that the resulting 244-state table does not implement the present
Routing245 source program.  The local theorem is a statement about all natural
register values and all finite lists of blocks.  Its proof below is by unary
and list induction.  The integration failures are separate, finite,
reproducible computations.  None of these results supplies a 244-state
consistency machine.

State numbers in this section are the original labels, with start state 95,
dispatcher 96, and external halt state 254.  The reference table has 245
working states and 490 transitions; these labels are distinct from the
canonical numbering with start 0, dispatcher 88, and halt 245.  The supplied
reference table is copied from the pinned Routing245 certificate, whose
provenance is recorded in the parent report and in the accompanying notices.

\subsection{Moving the register prefix instead of its suffix}

Write
\[
 B(v)=0\,1^{v+1},\qquad D(v)=1^{v+1}0
 \quad(v\in\mathbb N).
\]
For a finite list $A=(a_1,\ldots,a_k)$, let $B(A)$ and $D(A)$ be the
respective concatenations, and set
$w(A)=\sum_{i=1}^k(a_i+2)$.  Both concatenations are empty when $A$ is
empty.  A register selected for increment is marked as $00\,1^x$; its
unmarked representation would be $B(x)$.  The existing backend moves a
suffix and supplies additional spare space.  The experiment instead
slides the preceding register blocks and their header one cell left.

The experimental table makes exactly these three edits:
\begin{center}
\begin{tabular}{c|c|c}
State/read & Reference write/move/next & Experimental write/move/next\\
\hline
$75/0$  & $1,R,255$ & $0,L,112$\\
$83/0$  & $1,R,95$  & $1,L,112$\\
$131/0$ & $0,L,48$  & $0,R,48$
\end{tabular}
\end{center}
It then removes both transitions of state 255.  All other rows are
unchanged.  The file \texttt{left\_\allowbreak increment\_\allowbreak local244\_\allowbreak old.tm} has 244
working states and is retained as an explicitly rejected integration
experiment.

\paragraph{Word identity.}
For every finite natural-number list $A$,
\begin{equation}
 0D(A)=B(A)0.\label{eq:left-increment-word}
\end{equation}
The empty-list case is $0=0$.  Since $0D(v)=B(v)0$, the inductive step is
\[
 0D(v)D(A)=B(v)0D(A)=B(v)B(A)0.
\]
Thus this identity does not depend on testing finitely many block lengths.

\paragraph{One-block transitions.}
Starting in state 73 at the first symbol of $B(v)$, the machine passes the
block to the right unchanged, returning to state 73 just after the block,
in exactly $v+2$ moves.  The initial zero uses $73/0$ and the first one
uses $83/1$; subsequent ones use the state-73 self-loop.  The edited row
$83/0$ is consequently not used while seeking across an ordinary block.

Starting in state 131 at the rightmost symbol of $B(v)$, the machine
passes to its left, replacing $B(v)$ by $D(v)$, in exactly $v+2$ moves,
and again ends in state 131.  Indeed, $131/1$ writes zero and enters
state 112 moving left.  The $v$ remaining ones are traversed by the
$112/1$ self-loop.  The leading zero is replaced by one using $112/0$,
which returns to state 131 on its left.  This is a direct induction on
$v$.  Induction on the number of blocks then gives the same two scan
contracts for $B(A)$, with time $w(A)$, including $A=[]$.

\paragraph{Local increment theorem.}
Let $A,C$ be arbitrary finite lists of natural numbers and $x\in\mathbb N$.
Consider the finite tape segment
\begin{equation}
 \xi\;0\,0\,1\; B(A)\;0\,0\,1^x\;B(C)\;\zeta,
 \label{eq:left-increment-input}
\end{equation}
where the head is in state 73 at the first symbol of $B(A)$, or at the
first following marked zero if $A$ is empty.  The contexts $\xi,\zeta$
are arbitrary and will not be inspected.  After exactly $2w(A)+6$ moves,
the tape segment is
\begin{equation}
 \xi\;0\,1\; B(A)\;0\,1^{x+2}\;B(C)\;\zeta,
 \label{eq:left-increment-output}
\end{equation}
and state 48 scans the displayed header one immediately after $\xi0$.
The head has moved two cells left of its initial location.  The two
segments have equal length, so the unchanged right context occupies
exactly its original tape cells.

For the proof, first apply the state-73 list scan to $B(A)$, taking
$w(A)$ moves.  On the marked pair of zeros, the next three transitions are
\[
 (73,0)\mapsto(0,R,83),\quad
 (83,0)\mapsto(1,L,112),\quad
 (112,0)\mapsto(1,L,131).
\]
They leave $1^{x+2}$ to the right of the head.  The reverse state-131
list scan replaces $B(A)$ by $D(A)$ in $w(A)$ further moves.  At the old
header and its two preceding zeros, the final three transitions are
\[
 (131,1)\mapsto(0,L,112),\quad
 (112,0)\mapsto(1,L,131),\quad
 (131,0)\mapsto(0,R,48).
\]
The resulting word is
$\xi01\,0D(A)1^{x+2}B(C)\zeta$.
Applying~\eqref{eq:left-increment-word} gives
\eqref{eq:left-increment-output}.  The counts add to $2w(A)+6$;
none of the transitions enters $1^x$, $B(C)$, or either exterior context.
This also proves preservation of the suffix and the claimed uniformity
in $x$ and $C$.

The local effect is to increase the selected value by one and move its
prefix and header left by one cell.  The distance from the instruction
pointer to the header decreases by one.  The spare-block count is
unchanged.  In particular, the two zero cells immediately preceding the
header are a substantive precondition, not an optional formatting
convention.

\paragraph{Zero-decrement contraction.}
There is a second useful local fact.  When state 75 scans a zero whose
right neighbor is one, the reference machine performs
\[
 75/0\to255/1\to131/1
\]
in three moves, leaves both cells unchanged, and finishes in state 112
one cell to the left of the starting zero.  The replacement $75/0$
transition performs precisely that transformation in one move.  The
next-symbol premise is satisfied by the ordinary zero-decrement
register contract with a following nonempty block.  The corresponding
prefix carry and rollback scans yield time $2w(A)+5$, in place of
$2w(A)+7$.  This contraction is not valid at an increment tail whose
next cell is zero; the experiment changes that route jointly.

\texttt{prove\_left\_increment.py} checks the unary block derivations
and composes their list inductions using the inherited symbolic macro
checker.  It adds the explicitly proved identity
\eqref{eq:left-increment-word}.  Its output is labelled
\texttt{PASS\_LOCAL\_BOX\_CONTRACTS\_ONLY}.  This is a replayable
symbolic calculation with stated inference rules, accompanied here by
the mathematical induction proof; it is not a proof-kernel certification
of the whole experimental machine.

\subsection{The unchanged dispatcher does not supply the local premise}

During a genuine increment phase, the existing instruction protocol
fills the pointer-to-header gap with ones.  The register routine therefore
does not arrive at~\eqref{eq:left-increment-input}: the two required zero
cells are absent.  The attempted rollback crosses the header into the
instruction pointer.  A complete instruction fixture makes the failure
explicit.

Start at state 96 and head position zero, with pointer word $01010$,
header distance 181, physical primary registers
$(6,3,3,5,0,0,0,1,4)$, and 167 spare blocks of value one.  This is a
well-formed dispatcher fixture for logical instruction $\mathrm{CI}_8$,
whose physical register is 1.  After 976 moves, both tables enter state
73 at cell 182, immediately after the header at 181; cells 179 and 180
both contain one, directly exhibiting the failed local premise.
The edited machine reaches a clean
dispatcher frame after 1556 moves.  Its physical primary registers are
correctly $(6,4,3,5,0,0,0,1,4)$, but its pointer has ones at positions
$\{1,2\}$, represented by $0110$.  The required successor $01011$ has
ones at $\{1,3,4\}$.  The altered gap and spare count are also recorded:
\begin{center}
\begin{tabular}{l|r|c|r|r}
Execution & Moves & Pointer & Header distance & Spare count\\
\hline
Reference 245-state table & 2622 & $01011$ & 181 & 168\\
Edited 244-state table & 1556 & $0110$ & 180 & 167
\end{tabular}
\end{center}
Both outcomes are computed by independent sparse-set and dense-array
transition interpreters, which agree on every tape one at their stated
endpoints.  The changed pointer remains an error even if a proposed
simulation relation were to permit a different gap and spare count.

The fixture refutes the proposed unchanged-caller instruction contract.
It is not asserted to be reachable from the blank tape, and it does not
prove that every possible interpretation of the 244-state table is
incorrect.  Those stronger assertions are unnecessary: the proposed
modular proof cannot be composed, because this valid input to its
existing dispatcher contract has the wrong successor.

\subsection{A separate obstruction to automatic space closure}

Even a corrected dispatcher would need a space-allocation argument.
Suppose a proposed integration retains the present source execution and
its pointer-root-overflow events, and satisfies the following resource
rules: an increment decreases gap $g$ by one; a positive decrement
increases $g$ by one; each existing pointer-root overflow increases $g$
by one; other operations preserve $g$.  Let $I,D,O$ count these three
events since the seed, and let $S$ be the sum of the nine primary
registers.  The event counts are cumulative; an instruction may contribute
both a register event and an overflow event.

\paragraph{Resource invariant.}
For every finite source execution satisfying these rules,
\begin{equation}
 g=g_0+D-I+O,\qquad S=S_0+I-D,
 \qquad g+S-O=g_0+S_0.\label{eq:gap-invariant}
\end{equation}
The base case has all three counters zero.  An increment changes
$(g,S)$ by $(-1,+1)$, a positive decrement by $(+1,-1)$, and an overflow
changes $(g,O)$ by $(+1,+1)$.  Each preserves the last expression;
simultaneous events add their changes.  Operations of the remaining
types change none of these quantities.  Induction on the instruction
sequence proves all of~\eqref{eq:gap-invariant}.

For the actual Routing245 seed, $g_0=181$ and
the logical registers are $(3,0,0,0,3,0,0,0,0)$, so $S_0=6$ and
\begin{equation}
 g+S=187+O.\label{eq:seed-gap-invariant}
\end{equation}
Preserving the existing frame lower bound $g\ge64$ would require
$S\le123+O$ at every instruction boundary.

Two independent interpreters execute the exact 6460-leaf Routing245
source from this seed.  One follows a full binary trie and implements
carry and branch stripping directly on words.  The other indexes
zero-trimmed leaf words and uses integer addition and position lists.
They agree on every source transition and every event count for 3235
instructions.  Substitution into~\eqref{eq:seed-gap-invariant} gives:
\begin{center}
\begin{tabular}{r|r|r|r|r|r}
Source step & $I$ & $D$ & $O$ & $S$ & Implied $g$\\
\hline
2738 & 936 & 814 & 4 & 128 & 63\\
3235 & 1143 & 957 & 4 & 192 & $-1$
\end{tabular}
\end{center}
The first line is the first failure of $g\ge64$; the second is the first
negative gap.  The logical register vectors are respectively
$(7,6,0,0,4,0,34,77,0)$ and $(7,19,0,0,4,0,19,143,0)$.
These are exact source computations.  The gaps are counterfactual values
deduced for the proposed allocation policy, not measurements from a
verified 244-state execution.

Thus repairing only the local entry condition is insufficient for this
policy with the present source and control flow.  The negative gap also
rules out merely lowering the frame bound to zero.  An integration must
change the resource recurrence, or establish a changed source trace with
sufficient extra allocation events.  Grammar and control transformations
can change the frequency of root overflows; deliberate extension
operations change the recurrence.  The obstruction applies to the
specified policy and exact current execution, and does not establish a
lower bound for all Turing machines or all interpreter designs.

\subsection{Four reproducible finite compression screens}

The investigation also applied four bounded compression screens directly
to the exact 245-state table.  The supplied 8385 complete instruction
fixtures exercise all 490 transition rows of the reference table; the
longest reference fixture takes 2686 moves.  Each candidate run has fuel
20,000 and tape range $[-64,1983]$.  No candidate met all relevant
bounded checks:
\begin{center}
\begin{tabular}{l|r|r|r}
Screen & Trials & Past quick probes & Survivors\\
\hline
Selected-row state-pair merges, exact tape & 119560 & 54 & 0\\
Same pair merges, relaxed final frame & 119560 & 54 & 0\\
Single-entry state splices & 133712 & 163 & 0\\
Adaptive multi-entry splice search & 175680 & 712 & 0
\end{tabular}
\end{center}
For a pair merge, each of the two retained read-symbol rows is selected
from one of the two states, giving
$4\binom{245}{2}=119560$ candidates.  Single-entry splicing considers
137 states, excluding the start and dispatcher, with exactly one incoming
edge from outside the state.  It changes that edge to any write bit,
direction, and remaining target state: $137\cdot4\cdot244=133712$
trials.  Other transitions are preserved.

The multi-entry search first screens possible replacements for each
incoming edge independently against the otherwise unchanged machine.
It abandons a state as soon as an edge has no surviving replacement, and
then combines alternatives for covered states, with a cap of 100,000
combinations per state.  Its 175680 count is the number of individual
edge-edit trials, not a count of joint candidate tables.  Simultaneous
compensating edits that fail when applied alone are outside this search.

The relaxed endpoint allows a translated head, a header gap of at least
64, and any well-formed nonempty spare list.  It requires the same pointer
bits in the first 64 cells, the same nine primary values, and a clean
left marker.  This tests whether a local saving merely changes tape
resources.  The zero-survivor results hold even under this relaxation
for the stated search families.  They remain bounded computational
screens: fuel exhaustion is not a nonhalting proof, and failure on an
arbitrary valid fixture is not a proof about blank-tape reachability.

The four C++ sources, compressed fixture, recorded results, local symbolic
checks, and both exact counterexample scripts are included in
\texttt{interpreter/}.  They delimit several inexpensive dead ends and
isolate the missing ingredients of the more promising prefix-shifting
idea: a compatible instruction-pointer protocol and an unbounded
allocation mechanism with a proved frame invariant.
